%-------------------------------------- COMIENZA descripción del caso de uso.

  \begin{UseCase}{CU01}{Crear Mascota}{
      Permite al Propietario registrar una nueva mascota en el sistema, asociándola a su perfil para la gestión de su información y futuras reservaciones.
  }
      \UCitem{Versión}{\color{Gray}1.0}
      \UCitem{Autor}{\color{Gray}Castillo Aldaba Jared}
      \UCitem{Supervisa}{\color{Gray}Ugalde Tellez Aaron}
      \UCitem{Actor}{\hyperlink{Propietario}{Propietario}}
      \UCitem{Propósito}{Registrar los datos de una nueva mascota para gestionar su perfil, historial médico y futuras reservaciones.}
      \UCitem{Entradas}{
          \begin{itemize}
              \item \textbf{nombre}
              \item \textbf{especie\_id}
              \item \textbf{sexo\_id}
              \item \textbf{fecha\_nacimiento}
              \item Datos opcionales: raza\_id, peso\_kg, esterilizado, señas particulares, etc.
          \end{itemize}
      }
      \UCitem{Origen}{Teclado / Mouse}
      \UCitem{Salidas}{Mensaje de confirmación, Lista de mascotas actualizada.}
      \UCitem{Destino}{Pantalla y Base de Datos}
      \UCitem{Precondiciones}{
          \begin{itemize}
              \item El Propietario debe haber iniciado sesión en el sistema.
              \item El Propietario se encuentra en la pantalla de gestión de sus mascotas.
          \end{itemize}
      }
      \UCitem{Postcondiciones}{Se crea un nuevo registro en la tabla mascotas con los datos proporcionados, asociado al propietario\_id del usuario autenticado.}
      \UCitem{Errores}{
          \begin{itemize}
              \item \textbf{E1}: Datos de formulario inválidos o incompletos.
              \item \textbf{E2}: Usuario no autenticado o sesión expirada.
              \item \textbf{E3}: Error de conexión con el servidor.
          \end{itemize}
      }
      \UCitem{Tipo}{Primario}
  \end{UseCase}

  %--------------------------------------
  % Trayectoria Principal
  %--------------------------------------
  \begin{UCtrayectoria}
      \UCpaso[\UCactor] Accede a la \IUref{IU-PetList}{Lista de Mascotas} y oprime el botón \IUbutton{Agregar Mascota}.
      \UCpaso Despliega la \IUref{IU-PetForm}{Formulario de Mascota}.
      \UCpaso[\UCactor] Introduce los datos de la mascota en los campos correspondientes. \label{CU-CreatePet-FillForm}
      \UCpaso[\UCactor] Oprime el botón \IUbutton{Crear Mascota}.
      \UCpaso Verifica que los datos obligatorios (nombre, especie, sexo, fecha de nacimiento) estén completos y tengan el formato correcto con base en la regla \BRref{BR-PetValidation}{Validación de Datos de Mascota}.
      \UCpaso Solicita al servidor la creación del registro.
      \UCpaso El sistema verifica que el usuario esté autenticado y tenga permisos.
      \UCpaso El sistema registra la nueva mascota en la base de datos.
      \UCpaso Muestra el mensaje de éxito {\bf MSG1-}``Mascota creada exitosamente''.
      \UCpaso Actualiza la \IUref{IU-PetList}{Lista de Mascotas} para mostrar la nueva mascota.
  \end{UCtrayectoria}
  
  %--------------------------------------
  % Trayectorias Alternativas
  %--------------------------------------
  
  % Trayectoria A: Datos Inválidos
  \begin{UCtrayectoriaA}{A}{El Propietario ingresa datos inválidos}
      \UCpaso El sistema detecta que uno o más campos obligatorios están vacíos o tienen un formato incorrecto.
      \UCpaso Muestra el mensaje de error {\bf MSG-Error} ``Por favor completa todos los campos requeridos''.
      \UCpaso Resalta los campos con errores en el formulario.
      \UCpaso[\UCactor] Oprime el botón \IUbutton{Aceptar}.
      \UCpaso Continúa en el paso 5 de la trayectoria principal.
  \end{UCtrayectoriaA}
  
  % Trayectoria B: Sesión Expirada
  \begin{UCtrayectoriaA}{B}{La sesión del Propietario ha expirado}
      \UCpaso El sistema detecta que el token de autenticación es inválido o ha expirado.
      \UCpaso Muestra el mensaje de error {\bf MSG-Error} ``Tu sesión ha expirado. Por favor, inicia sesión de nuevo.''.
      \UCpaso Redirige al usuario a la \IUref{IU-Login}{Pantalla de Inicio de Sesión}.
      \UCpaso Termina el caso de uso.
  \end{UCtrayectoriaA}
  
  % Trayectoria C: Fallo del Sistema
  \begin{UCtrayectoriaA}{C}{Fallo inesperado del sistema}
      \UCpaso Ocurre un error interno en el servidor al intentar registrar la mascota.
      \UCpaso Muestra el mensaje de error {\bf MSG-Error} ``Ocurrió un error al procesar tu solicitud. Por favor, intenta más tarde.''.
      \UCpaso[\UCactor] Oprime el botón \IUbutton{Aceptar}.
      \UCpaso[] El sistema permanece en el \IUref{IU-PetForm}{Formulario de Mascota}, permitiendo al usuario reintentar la operación.
  \end{UCtrayectoriaA}
%--------------------------------------
% Puntos de extensión
%--------------------------------------
\subsection{Puntos de extensión}
\UCExtenssionPoint{
	% Cuando:
	El propietario desea ver el perfil completo y agregar más detalles (vacunas, enfermedades, etc.) a la mascota recién creada.
}{
	% Durante la región:
	Paso 10 (Visualización de la lista actualizada).
}{
	% Casos de uso a los que extiende:
	\hyperlink{CU14}{CU14 Consultar Detalle de Mascota}, \hyperlink{CU06}{CU06 Agregar Historial Médico}.
}

\UCExtenssionPoint{
	% Cuando:
	El propietario desea crear una reservación para la mascota recién creada.
}{
	% Durante la región:
	Paso 10 (Visualización de la lista actualizada).
}{
	% Casos de uso a los que extiende:
	\hyperlink{CU15}{CU15 Crear Reservación}.
}

%-------------------------------------- TERMINA descripción del caso de uso.

